\documentclass[reqno,11pt]{article}
\usepackage[
backend=biber,
style=numeric,
sorting=ydnt,
doi={false}
]{biblatex}

\addbibresource{savare.bib}



\usepackage{enumitem}
\usepackage{eurosym}  
\usepackage{multirow}
\usepackage{xcolor}
\usepackage{colortbl}
%\usepackage{parskip}
 \usepackage{newpxtext,newpxmath}
%\linespread{1.08}
\setlength{\parskip}{1.5pt plus 0.5pt minus 0.5pt}
\let\OLDthebibliography\thebibliography
\renewcommand\thebibliography[1]{
  \OLDthebibliography{#1}
  \setlength{\parskip}{0pt}
  \setlength{\itemsep}{0pt plus 0.3ex}}
\setlength{\parindent}{0pt}


\newenvironment{myquote}{\list{}{\leftmargin=0.25in\rightmargin=0.25in}\item[]}{\endlist}
%\renewcommand\thesection{\Alph{section}}
\newenvironment{smallquote}{\list{}{\leftmargin=0.15in\rightmargin=0.15in\topsep=-4pt}\item[]}{\endlist}

\usepackage{titlesec}
\titleformat*{\subsection}{\large\bfseries}
\titlespacing*{\subsection}{0pt}{1ex plus 1ex minus .2ex}{0.5ex plus
  .2ex}
\titlespacing*{\section}{0pt}{1ex plus 1ex minus .2ex}{0.5ex plus
   .2ex}

 \titleformat*{\subsubsection}{\bfseries}{}{}{}{}
\titlespacing*{\subsubsection}{0pt}{0.5ex plus .2ex minus .2ex}{0.5ex plus
   .2ex}
 \let\mysec\subsubsection

 \newcommand{\present}{}


\AtBeginDocument{
\addtolength{\abovedisplayskip}{-1.5ex}
\addtolength{\abovedisplayshortskip}{-.7ex}
\addtolength{\belowdisplayskip}{-1.5ex}
\addtolength{\belowdisplayshortskip}{-.7ex}
}

%\setlength{\parindent}{0pt}
\newcommand{\pind}{\hangindent \parindent
  \indent} 
\newcommand{\npind}{\hangindent\parindent
  \noindent} 

%\newcommand{\problem}[1]{\npind\fbox{\footnotesize#1}\quad}

\newcommand{\point}[1]{\fbox{\scriptsize #1}}
\newcommand{\squig}{{\ensuremath\rightsquigarrow}}
\newcommand{\problem}[3]{%\smallskip
  \npind
  \fbox{\scriptsize #1}
  \hspace{3pt}\textbf{#2}\\#3}

\newenvironment{tightcenter}{%
  \setlength\topsep{2pt}
  \setlength\parskip{0pt}
  \begin{center}
}{%
  \end{center}
}

\newenvironment{reducedcenter}{%
  \setlength\topsep{2pt}
  \setlength\parskip{0pt}
  \begin{center}
}{%
  \end{center}
\vspace{-4pt}
}
\newcommand{\myskip}{\smallskip\noindent}

\newcommand{\workshop}[4]{#1, #2, #3, {\em #4}}


%\bibliography{bibliografia2014.bib}
\usepackage{fancyhdr}

%\rhead{Curriculum}
%\lhead{Giuseppe Savar\'e}
%\chead{}
%headsep=12pt,
\usepackage[left=2.5cm, right=2.5cm,top=0.6cm,bottom=0.6cm,footskip=18pt,
includehead, includefoot,headsep=12pt]{geometry}

\renewcommand{\thesubsection}{\arabic{subsection}}
\let\mysec\subsubsection
\newcommand{\switch}[2]{\iflanguage{italian}{#1}{#2}}

\pagestyle{fancy}
\fancyhf{}
\renewcommand{\headrulewidth}{0pt}
%\lhead{\small\itshape Giuseppe Savar\'e -- Curriculum}
\rhead{\small\thepage}


\newcommand{\stateoftheart}{%
%\vspace{-12pt}
\begin{center}
  \bf \large State of the art
\end{center}}

\newcommand{\obmeth}{%
  \begin{reducedcenter}
    \bf \large Objectives and methodology.
  \end{reducedcenter}}
\newcommand{\objectives}{%
  \begin{center}
    \bf \large Objectives.
  \end{center}}
\newcommand{\meth}[1]{\smallskip\noindent
  \kern-39pt
  \fbox{$
    {\text{\scriptsize
        %\tiny
        \emph{%
          %\kern1.5pt 
          Methodology%
          %\kern1.5pt
        }}} 
  $}
\hspace{3pt}\npind
%\begin{minipage}[t]{0.85\linewidth}
    #1
%  \end{minipage}
}

\newcommand{\timescale}[1]
{\smallskip\noindent
  %\kern-48pt
  \fbox{$
    \stackrel{\text{%
        \scriptsize
        %\tiny
        \emph{Timescale and}}} {{\text{%
         \scriptsize
         % \tiny
         \emph{resources}}}}$}
  \hspace{6pt}
  %\begin{minipage}[t]{0.85\linewidth}
    #1
%  \end{minipage}
  }
\newcommand{\team}[1]{  % 
  \begin{flushright}
    \begin{minipage}{0.8\linewidth}
      % \textrm{Team:} 
      \em #1
      \end{minipage}
\end{flushright}}


%  \newcommand{\team}[1]{}
%{  % 
%\vspace{-8pt}
%\begin{flushright}
%\textrm{Team:} \em #1
%\end{flushright}}

\usepackage{eucal,enumerate,mathrsfs,xspace}
\let\Bbbk\relax
\usepackage{amsmath,amssymb,epsfig,theorem}
%\usepackage[toc,page]{appendix}

\renewcommand{\phi}{\mathscr E}

\usepackage{ifthen}
%\numberwithin{equation}{section}


\theoremstyle{plain}
\newtheorem{maintheorem}{Theorem}
\newtheorem{maincorollary}{Corollary}

\newtheorem{theorem}{Theorem}[section]
\newtheorem{corollary}[theorem]{Corollary}
%\newtheorem{problem}[theorem]{Problem}
\newtheorem{lemma}[theorem]{Lemma}
\newtheorem{proposition}[theorem]{Proposition}
\newtheorem{definition}[theorem]{Definition}
\newtheorem{notation}[theorem]{Notation}
\newtheorem{hypothesis}[theorem]{Hypothesis}
\newtheorem{assumption}[theorem]{Assumption}
%\newtheorem*{definition*}{Definition}


%\theoremstyle{remark}
\newtheorem{remark}[theorem]{Remark}
\newtheorem{example}[theorem]{Example}
\newtheorem{claim}{Claim}
%\newtheorem*{claim*}{Claim}
%\newtheorem*{remark*}{Remark}
%\newtheorem*{example*}{Example}
%\newtheorem*{notation*}{Notation}



%%%%%%%%%%%%%%%%%%%%%%%%%%%%%%%%%%%%%%%%%%%%%%%%%%%%%%%%%%%
%                                                        % Packages for drafts, labels, corrections
%%%%%%%%%%%%%%%%%%%%%%%%%%%%%%%%%%%%%%%%%%%%%%%%%%%%%%%%%%%

%\usepackage[normalem]{ulem}
%\usepackage[color]{showkeys}
\usepackage{pdfsync}
\usepackage[colorlinks,linkcolor=blue,citecolor=blue,urlcolor=blue]{hyperref}
\newcommand{\plambda}{$\lambda$--}
 
%%% Macros (inlined from macro2014.tex and macro_metric2014.tex)
\newcommand{\V}{\mathbb{V}}
\newcommand{\cE}{{\ensuremath{\mathcal E}}}
\newcommand{\sfS}{{\sf S}}
\newcommand{\sfW}{{\sf W}}
\newcommand{\sfd}{{\sf d}}
\newcommand{\Cheeger}{\mathsf{Ch}}
\newcommand{\Ch}{\Cheeger}
%
%
%%%%%%%%%%%%%%%%%%%%%%%%%%%%%%%%%%%%%%%%%%%%%%%%%%%%%%%%%%%%
%                                                         % MACRO SPECIFICHE
%%%%%%%%%%%%%%%%%%%%%%%%%%%%%%%%%%%%%%%%%%%%%%%%%%%%%%%%%%%%
%
%
\newcommand{\sfde}{\ssfd_\cE}
\newcommand{\one}{{{\bf 1}}}
\newcommand{\Vrestr}[1]{\V\kern-2pt_{#1}}
\newcommand{\Cp}{\mathrm{Cap}}
\newcommand{\bemph}[1]{\textbf{\emph{#1}}}

%
%

\newcommand{\ssfd}{\scalebox{0.95}{\sfd}}
\newcommand{\ssfS}{\scalebox{0.9}{\sfS}}
\newcommand{\ssfW}{\scalebox{0.9}{\sfW}}

\renewcommand{\Ch}{\scalebox{0.95}{$\mathsf{Ch}$}}

\newcommand{\HK}{\mathsf{H\kern-3pt K}} 


%\newcommand{\mymathsf}[1]{\mathsf{#1}}

%

\newcommand{\mybreak}{\smallskip\noindent}
\newcommand{\biotitle}[1]{\mybreak #1\par\mybreak}




\begin{document}

\thispagestyle{empty}
% \nocite{*}
% \maketitle

\setcounter{page}{1}
%\cfoot{B1--b\thepage}
\titleformat{\subsubsection}{\bfseries}{}{0pt}{}[\vspace{0.3ex}\hrule]
\titlespacing*{\subsubsection}{0pt}{.8ex plus .2ex minus .2ex}{.4ex plus
   .2ex}

\begin{center}
    \Large Giuseppe Savar\'e -
    CURRICULUM
\end{center}


\mysec*{PERSONAL INFORMATION}

\begin{tabular}{lll}
Savar\'e Giuseppe&
ORCID:
  \href{https://orcid.org/0000-0002-0104-4158}{0000-0002-0104-4158}
  &MathSciNet:	\href{https://mathscinet.ams.org/mathscinet/MRAuthorID/336952}{336952}\\
ResearcherID: \href{https://www.webofscience.com/wos/author/record/H-3651-2014}{H-3651-2014}
&Scopus: \href{https://www.scopus.com/authid/detail.uri?authorId=6602154611}{6602154611}
&Google Scholar: \href{https://scholar.google.com/citations?user=w7qH1YYAAAAJ&hl=en}{w7qH1YYAAAAJ}\\
Nationality: Italian &
Born in Pavia, 25/08/1966\\
\multicolumn{3}{l}{Homepage: \href{https://giuseppesavare.com}{giuseppesavare.com} \quad Email: \href{mailto:giuseppe.savare@unibocconi.it}{giuseppe.savare@unibocconi.it}}\\
\multicolumn{3}{l}{Department of Decision Sciences, Bocconi University, Via Roentgen 1, 20136 Milano, Italy}
\end{tabular}

\mysec*{RESEARCH INTERESTS}
Optimal Transport and entropic distances between nonnegative measures;
gradient flows, rate-independent problems, and applications to evolution PDEs;
analysis in metric-measure spaces: calculus, heat flow, optimal transport, non-smooth Riemannian geometry;
variational methods for Partial Differential Equations.

\mysec*{EDUCATION}

%\mybreak
1984--1988: {\em Laurea (cum laude) in Mathematics},
University of Pavia, Italy
%\pind{\em Problemi
%di discretizzazione per equazioni differenziali astratte}.
%\switch{Relatore:}{Advisor:} prof.~Claudio Baiocchi.

\mysec*{CURRENT POSITION}
2020--: \emph{Full Professor,} Mathematical Analysis. 
Department of Decision Sciences, Bocconi University

\mysec*{PREVIOUS POSITIONS}
2000--2020: \switch{Professore ordinario di Analisi
  Matematica}{\emph{Full Professor,}
  Mathematical
  Analysis}
Department of Mathematics,   University of
Pavia

\noindent
1998--2000: 
\switch{\emph{Professore associato} di Analisi Matematica}{\emph{Associate professor,}}
%Mathematical Analysis}
University of Pavia

\noindent
1990--1998: \emph{\switch{Ricercatore}{Researcher,}}
%\pind
\switch{Istituto di Analisi Numerica del C.N.R.}{IMATI-CNR, Pavia}



\mysec*{FELLOWSHIPS AND AWARDS}

\hangindent \parindent
\noindent 
2025: ERC Advanced Grant, awarded for the project
\href{https://dec.unibocconi.eu/news/giuseppe-savare-awarded-erc-advanced-grant}{\emph{Optimal Transport and Metric Structures
for Evolution Problems}} (OPTiMiSE)

\npind
2025: Included in the \href{https://clarivate.com/highly-cited-researchers/}{\emph{Highly Cited Researchers}} list by Clarivate for the field of Mathematics

\npind
2019--2023:
\href{https://www.ias.tum.de/ias/programs/fellowship-programs/hans-fischer-senior-fellowship/}{\emph{Hans Fischer senior fellow}}, Institute for Advanced Study,
Technical University of Munich

\npind
2019: \emph{Plenary Speaker}, XXI Congresso Unione Matematica Italiana

\hangindent \parindent
\noindent 
2016: \emph{Invited Speaker}, 7$^{th}$ \href{https://www.7ecm.de/home/}{European Congress of Mathematics},
Berlin: section ``Analysis and PDEs''

\hangindent \parindent
\noindent 
2015: \emph{John von Neumann Visiting Professor},
Technical University of Munich

\hangindent \parindent
\noindent 
2014: the results on Gradient Flows have been presented in the
\href{https://smf.emath.fr/publications/expose-bourbaki-1065-flots-de-gradient-dans-les-espaces-metriques-et-leurs}{\emph{Expos\'e Bourbaki 1065: Flots de gradient dans les espaces
  m\'etriques et leurs applications d'apr\`e Ambrosio-Gigli-Savar\'e}},
Asterisque 361.

\hangindent \parindent
\noindent 
  2012: The paper \emph{Chemical reactions as $\Gamma$-limit of
    diffusion}
  SIAM Rev. 54 (2012) (with M.~Peletier, M.~Veneroni) 
  has been published as SIAM Review’s SIGEST award.

\hangindent \parindent
\noindent 2011: 
\href{https://umi.dm.unibo.it/premi-old/premio-ennio-de-giorgi/}{\emph{Ennio De Giorgi Prize}},
  awarded by the
\emph{Italian Mathematical Union}
% \switch{per l'Analisi Matematica}{to a 
% mathematician under 45 years of age}

% \hangindent \parindent
% \noindent 
% 2009--present: Membro Corrispondente, \emph{Istituto
%   Lombardo, Accademia di Scienze e Lettere}, Milano.


\hangindent \parindent
\noindent
1994: 
\emph{Gioachino Japichino} Prize, awarded by
the
\emph{Accademia Nazionale dei Lincei} 
%to a mathematician under 30 years of age


% \mysec*{FELLOWSHIPS AND AWARDS}

% \hangindent \parindent
% 1994: 
% \emph{Gioachino Japichino} Prize, awarded by
% the
% \emph{Accademia Nazionale dei Lincei} 
% to a mathematician (under 30 years of age) for a relevant publication
% in Analysis.

% \hangindent \parindent
% \noindent 2011: 
% \emph{Ennio De Giorgi Prize,} first edition;
%   awarded by the
% \emph{Italian Mathematical Union}
% \switch{per l'Analisi Matematica.}{to a 
%   mathematician under 45 years of age.}


% \hangindent \parindent
% \noindent 
% 2015: \emph{John von Neumann Visiting Professor},
% Technische-Universit\"at M\"unchen.

% \hangindent \parindent
% \noindent 
% 2016: Invited Speaker at the 7$^{th}$ European Congress of Mathematics
% (7ECM), Berlin: section ``Analysis and PDEs''.


\mysec*{SUPERVISION OF GRADUATE STUDENTS}

% \npind
% 5 PhD thesis: Simona Sanfelici 1998, 
% Riccarda Rossi 2005, Stefano Lisini 2006, Marco Veneroni 2007
% (now associate or assistant professors), and Luca Natile, 2011
% (now high school professor).

\hangindent \parindent
\noindent \emph{Bocconi University}:\\
current: Alessandro Pinzi

\hangindent \parindent
\noindent \emph{IAS-TUM, Munich}:\\
2022: \href{https://sites.google.com/view/giacomoenricosodini}{Giacomo
Sodini}
(co-advised with M.~Fornasier,
now PostDoc, TU Wien)

\hangindent \parindent
\noindent \emph{University of Pavia}:\\
2019:
\href{https://sites.google.com/view/nicolodeponti}{Nicol\`o
De Ponti}
(now Tenure-Track Researcher, Politecnico di Milano)\\
2017:
Luca
Minotti\\
2015:
Giovanni
Bonaschi (co-advised with M.~Peletier, Eindhoven; now Portfolio manager fondi Multi Asset)\\
2014:
\href{https://dariomazzoleni.weebly.com}{Dario
Mazzoleni} (co-advised with A.~Pratelli, Erlangen;
now Associate Professor, University of Pavia)\\
2011:
Luca
Natile\\
2007:
\href{https://mate.unipv.it/~veneroni/}{Marco
Veneroni}
(now Associate Professor, University of Pavia)\\
2006:
\href{https://mate.unipv.it/lisini/}{Stefano
Lisini}
(now Associate Professor, University of Pavia)\\
2005:
\href{https://riccarda-rossi.unibs.it}{Riccarda
Rossi}
(now Associate Professor, University of Brescia)\\
1998:
\href{https://personale.unipr.it/it/ugovdocenti/person/19220}{Simona
Sanfelici} (co-advised with P.~Colli Franzone;
now Full Professor, University of Parma)\\
%
%
%
%
%
\smallskip\noindent
I also mention the supervision of the Master thesis
(subsequently published in international mathematical journals) of
\href{https://sites.google.com/view/federicobassetti/}{Federico Bassetti} (now Full Professor, Politecnico di Milano),
\href{http://arturo.imati.cnr.it/spinolo/}{Laura Spinolo} (now Senior Researcher, IMATI-CNR),
\href{https://www.gssi.it/people/professors/lectures-maths/item/2297-daneri-sara}{Sara Daneri} (now Associate Professor, GSSI L'Aquila),
\href{https://giulslu.github.io}{Giulia Luise} (now Researcher, Microsoft)


\mysec*{POSTDOCTORAL FELLOWS}

2024--2025: \href{https://sites.google.com/view/filipporiva/home/}{Filippo Riva}
(now PostDoc, Charles University, Prague)\\
2021--2023:
\href{https://sites.google.com/view/lucatamanini}{Luca
Tamanini}
(now Assistant Professor, Universit\`a Cattolica del Sacro Cuore, Brescia)\\
2017--2019:
\href{https://sites.google.com/view/giulia-cavagnari}{Giulia
Cavagnari}
(now Associate Professor, Politecnico di Milano)\\
2016--2017:
\href{https://dariomazzoleni.weebly.com}{Dario
Mazzoleni}
(now Associate Professor, University of Pavia)\\
2015--2016:
\href{https://mate.unipv.it/orrieri/}{Carlo
Orrieri}
(now Associate Professor, University of Pavia),
\href{https://www4.ceda.polimi.it/manifesti/manifesti/controller/ricerche/RicercaPerDocentiPublic.do?evn_prodotti=EVENTO&k_doc=265344&aa=2021&lang=IT}{Matteo
Muratori}
(now Full Professor, Politecnico di Milano)\\
2009--2011:
\href{http://test.dime.unige.it/it/users/edoardo-mainini}{Edoardo
Mainini}
(now Associate Professor, University of Genova)\\
2007--2008:
\href{https://www.professoren.tum.de/matthes-daniel}{Daniel
Matthes}
(now Associate Professor, TUM)\\
2006--2008:
\href{https://www.di.univr.it/?ent=persona&id=6498&lang=en}{Antonio
Marigonda}
(now Full Professor, University of Verona)\\

% \npind 
% 2006-2008: Antonio Marigonda (now assistant professor,
% University of Verona) 

% \npind 
% 2007-2008: Daniel Matthes (now full professor, Technische
% Universit\"at M\"unchen)

% \npind
% 2009-2011: Edoardo Mainini (now assistant professor, University of Genova)
% % \npind 



% % \npind 
% % 2002: Laura Spinolo (now researcher, IMATI-CNR)

% % \npind
% % 2002: Luca Heltai (now assistant professor, SISSA, Trieste)

% % \npind
% % 2007: Sara Daneri (now Research fellow, University of Leipzig)

\mysec*{TEACHING ACTIVITIES}
\pind
\emph{University of Pavia}:\\
1998--2020: courses of \emph{Mathematical Methods}
(Engineering Faculty) and \emph{Mathematical
  Analysis} (degree in Mathematics), typically two main courses
every year,$\sim$140 hours/year\\
2000--2020:
several Ph.D. courses (PDEs, Semigroups, Calculus of Variations, Optimal transport)

\npind
\emph{Bocconi University}:\\
2020--2025: Mathematical Analysis (Bachelor), 
Introduction to Real Analysis (Ph.D.)
\\
2025: Alegbraic and Topological Methods (Bachelor)

\npind
Several invited courses to \emph{international advanced schools}
(as CIME, HIM, SNS, EVEQ)
% \mysec*{TEACHING ACTIVITIES}
% \pind
% I have mainly been teaching courses of %Advanced Calculus 
% Mathematical Methods for the Engineering Faculty.

% \noindent
% Since 2000, for the PhD programme in \emph{Mathematics and
%   Statistics} (Pavia) I have also been teaching various courses as 
% %Mathematical models in Electrocardiology. 
% PDEs,
% Semigroup Theory, % Evolution equations in
% % Banach spaces. 
% Calculus of Variations, Optimal transport.
% % Transport problems,
% % Wasserstein distance and evolution equations.
% %Optimal transport, diffusion equations and Riemannian geometry.
% %Variational methods for evolution problems.
% % \begin{tabular}{ll}
% % Mathematical models in Electrocardiology, &Partial
% % differential equations, \\
% % Semigroup Theory, & Evolution equations in
% % Banach spaces,\\
% % Calculus of Variations, &Optimal transport, \\
% % \multicolumn{2}{l}{Transport problems,
% % Wasserstein distance and evolution equations,}\\
% % \multicolumn{2}{l}{Optimal transport, diffusion equations and Riemannian geometry,}\\
% % Variational methods for evolution problems.
% % \end{tabular}
% \\
% \noindent A list of invited courses to international advanced schools is
% included in the next section c.

% \switch{Si \`e principalmente svolta nell'ambito dei corsi:}{My teaching activity has been mainly concerned with the
%  following courses:}
% \begin{itemize}
% \item{\sl Analisi Matematica I} \switch{}{(basic calculus for
%   functions of one variable)}
% \item{\sl Analisi Matematica
%     II} \switch{}{ (advanced calculus)}
% \item{\sl Metodi Matematici
%     per l'Ingegneria} \switch{(Laurea triennale e laurea magistrale)}{(Mathematical Methods for Engineering)}

  
% \item{\sl Modelli Matematici per l'elettrocardiologia} 
%   \switch{(dottorato)}{ (Mathematical models in Electrocardiology, Ph. D Course)}
% \item{\sl Equazioni alle derivate parziali} 
% \switch{(dottorato)}{(Partial
%   differential equations, Ph. D. course)}
% \item{\sl Teoria
%     dei semigruppi} \switch{(dottorato)}{(Semigroup theory, Ph. D. course)}
  
%   \item{\sl Equazioni di evoluzione negli spazi di
%       Banach} \switch{(dottorato)}{(Evolution equations in Banach spaces, Ph. D course)}
%   \item{\sl Problemi di trasporto, distanza di
%     Wasserstein ed equazioni di evoluzione} 
%   \switch{(dottorato)}{(Transport problems,
%   Wasserstein distance and evolution equations, Ph. D. course)}
% \end{itemize}



\mysec*{ORGANISATION OF SCIENTIFIC MEETINGS}
2026: MFO Workshop
\href{https://www.mfo.de/www/activity/2613}{\emph{Flows on Measure Spaces and Applications in Machine Learning}},
Oberwolfach
{\small (Co-organizers: P.~Rigollet, G.~Steidl, F.-X.~Vialard)}

2026: Intensive school \href{https://sites.google.com/unitn.it/mfml-2026/home}{\emph{Mathematical Foundations of Machine Learning (MFML 2026)}},
%May 25-29,
Vason (Monte Bondone), Trento.

2025: International School 
\href{https://mmaa.lakecomoschool.org/}{\emph{Mathematical Analysis and Applications}}, 
%June 9-13 2025, 
Lake Como School of Advanced Studies, Villa del Grumello, Como  
{\small (Co-organizer: M.G.~Mora)}

2025: Workshop on \href{https://bidsa.unibocconi.eu/variational-methods-and-non-smooth-geometric-structures-applications}{\emph{Variational Methods and Non-Smooth Geometric Structures with Applications}}, June 3-5 2025, Bocconi University, Milan.

2024: Workshop on \href{https://dec.unibocconi.eu/mathematics-artificial-intelligence-and-machine-learning}{\emph{Mathematics for Artificial Intelligence and Machine Learning}}, Universit\`a Bocconi, Milano (Co-organizers: B. Morini, E. De Vito, S. Pieraccini)

2023: Workshop on \href{https://sites.google.com/view/vargeo2023/}{\emph{Variational and geometric structures for evolution}}, CIRM, Trento
{\small (Co-organizers: D. Knees, R. Rossi, M. Thomas)}

2023: \href{https://www.ias.tum.de/ias/event-pages/workshop-otmfml/home/}{\emph{Workshop on Optimal Transport, Mean-Field Models, and
  Machine Learning}}, IAS-TUM, Munich
{\small (Co-organizers: M. Burger, M. Fornasier,
  G. Peyr\'e)}
  
\npind
2026, 2024, 2022, 2018, 2016, 2014, 2012, 2010, 2008: \href{http://www.crm.sns.it/event/501/}{\emph{Workshops on Optimal
Transportation and Applications}},
Pisa {\small (Co-organizers: L. Ambrosio, G. Buttazzo, N. Gigli)}

\npind
2022: \href{https://contrendspde.wordpress.com}{\em Contemporary Trends in Kinetic Theory and PDEs}, Pavia {\small (Co-organizers: J.A. Carrillo, A. Pulvirenti, M. Zanella)}

\npind
2018: School on \href{https://otnm.lakecomoschool.org}{\emph{Optimal transport: numerical methods and
  applications}},
Lake Como School of Advanced Studies {\small (Co-organizer: F.~Santambrogio}).

\npind
2018: Workshop \href{https://sites.google.com/view/controlpv2018}{\emph{Optimal Control and Mean Field Games}}, Pavia
{\small (Co-organizers: G. Cavagnari, S. Lisini, C. Orrieri)}


\npind
2016: Bimester on \href{https://appliedmath.univie.ac.at/public/ESI_nonlinear_flows/}{\emph{Nonlinear Flows}}, Research Centre ESI,
Vienna
{\small (Co-organizers: E. Feireisl, A. Juengel, A. Mielke, U. Stefanelli)}

\npind
2014 and 2011: MFO Workshop \emph{Variational Methods for Evolution},
Oberwolfach {\small (Co-organizers: L. Ambrosio, A. Mielke,
  M. Peletier, F.~Otto, U. Stefanelli)}

\npind
2011: Conference \emph{Analysis and Numerics of PDEs. In memory of Enrico Magenes}, Pavia

\npind
2010: BIRS Workshop: \emph{Rate-independent systems: Modeling,
  Analysis, and Computations},
Banff {\small (Co-organizer: U. Stefanelli)}

\npind
2008: CIME Course: \emph{Nonlinear Partial Differential Equations and
  Applications}, Cetraro {\small (Co-organizer: L. Ambrosio)}



% \pind 2016: bimester on \emph{Nonlinear Flows}, 
% Research Centre ESI, 
% University of Vienna.\\
% Co-organizers:
% E.~Feireisl, A.~Juengel, A.~Mielke, U.~Stefanelli.

% \npind
% 2011 and 2014: MFO Workshop.
% \emph{Variational Methods for Evolution},
% Oberwolfach (DE).\\
% Co-organizer: L. Ambrosio, A. Mielke, M. Peletier, U. Stefanelli.


% \npind
% 2008, 2010, 2012, 2014: Workshops on \emph{Optimal Transportation and Applications},
% De Giorgi Center, Pisa (IT).
% Co-organizer: L. Ambrosio, G. Buttazzo.

% \npind
% 2011: Conference 
% \emph{Analysis and Numerics of PDEs - In memory of Enrico Magenes,} 
% Pavia (IT).

% \npind
% 2010: BIRS Workshop: \emph{Rate-independent systems: Modeling, Analysis, and
% Computations}, Banff (CA).
% Co-organizer: U. Stefanelli.


% \npind 
% 2008: CIME Course: \emph{Nonlinear Partial Differential Equations and Applications,}
% Cetraro (IT).\\
% Co-organizer: L. Ambrosio.



\mysec*{INSTITUTIONAL RESPONSIBILITIES}
2026--\present: Head, Department of Decision Sciences, Bocconi University\\
\pind
2019: Director of the \emph{Advanced School of Ph.D.~Higher Education}
(SAFD), University of Pavia\\
2014--2016:
\switch{Membro del Nucleo di Valutazione,}{\emph{University Assessment
    Commission}}, University of Pavia\\
2001--2008: Director of the \emph{Ph.D.~program in Mathematics and
  Statistics}, University of Pavia\\
\npind
1998--\present: IMATI-CNR, research associate




\mysec*{EDITORIAL AND ADVISORY BOARDS}


\pind
2023--\present: \href{https://www.esaim-cocv.org}{ESAIM: Control, Optimisation and Calculus of Variations},
member of the \emph{editorial board}\\
2018–\present: Unione Matematica Italiana, member of the \emph{Scientific Advisory Board} % (UMI).  

\npind
2014--\present:
C.I.M.E.~Foundation, member of the \emph{Scientific Advisory Board}
% \switch{Membro del Consiglio Scientifico della fondazione
%   C.I.M.E.}{\emph{Scientific Advisory Board}, C.I.M.E.~foundation}

\npind
2016–\present: \href{https://www.springer.com/journal/245}{Applied Mathematics and Optimization}, member of
the \emph{editorial board}\\
2013–\present: \href{https://www.springer.com/journal/11118}{Potential Analysis}, member of the \emph{editorial board}

\mysec*{MEMBERSHIPS OF SCIENTIFIC SOCIETIES}

\hangindent \parindent
\noindent 
2009--\present: Membro Corrispondente, \emph{Istituto
  Lombardo, Accademia di Scienze e Lettere}, Milano

% \npind
% 1998–\present: IMATI-CNR, research associate  



% \mysec*{COMMISSIONS OF TRUST}


% 2014--present:
% \switch{Membro del Consiglio Scientifico della fondazione
%   C.I.M.E.}{Scientific Advisory Board, C.I.M.E.~foundation, Firenze, Italy.}

% \npind
% 2014--\present:
% \switch{Membro del Nucleo di Valutazione,}{University Assessment Commission,}
% University of Pavia.


%\mysec*{MEMBERSHIPS OF SCIENTIFIC SOCIETIES}


% 201? –	Member, Research Network “Name of Research Network”
% 200? –	Associated Member, Name of Faculty/ Department/Centre, Name of University/ Institution/ Country
% 200? –	Funding Member, Name of Faculty/ Department/Centre, Name of University/ Institution/ Country 



% \mysec*{MAJOR COLLABORATIONS}
% I have collaborated with more than 50 different co-authors. Among the main
% collaborations (not counting former Ph.D.~students and post-docs),
% I recall



% \npind 
% \emph{Luigi Ambrosio} (Scuola Normale Superiore, Pisa, I): %since 1998:
% Optimal Transport, Gradient flows, analysis in metric spaces.
% Together with N. Gigli, we wrote the monograph
% \emph{Gradient flows in metric spaces and in spaces of probability
%   measures,}
% Birkh\"auser, 2005 (second edition in 2008)


% \npind
% \emph{Piero Colli Franzone} (Pavia): %since 1996: 
%  mathematical models for electrocardiology



%  \npind
% \emph{Jean Dolbeault} (Universit\'e Paris Dauphine, Paris): %since 2006:
% Transport distances, functional inequalities

% \npind
% \emph{Massimo Fornasier} (Technical University of Munich): mean-field
% optimal control


% \npind
% \emph{Alexander Mielke} (WIAS, Berlin): %since 2005:
% Rate-independent processes, doubly-nonlinear evolution equations,
% Optimal Entropy-Transport problems


% \npind
% \emph{Nicola Gigli} (SISSA, Trieste):
% %since 2002:
% % Optimal Transport, Gradient flows, 
% analysis in metric-measure spaces

% % \npind
% % \emph{Robert McCann} (Toronto, CA): Fourth order nonlinear diffusion
% % equations

% \npind
% \emph{Ricardo H. Nochetto} (University of Maryland, USA): %1996-2006:
% optimal error estimates for %nonlinear 
% evolution problems

% \npind
% \emph{Mark Peletier} (TU
% Eindhoven, NL): %since 2008:
% Reaction-diffusion systems, Fokker-Planck equations

% \npind
% \emph{Giuseppe Toscani} (Pavia, I): %since 2006:
% nonlinear diffusion equations, entropy methods. % functional inequalities.


% I also collaborated with \emph{Yann Brenier, Jos\'e A. Carrillo,
%   Robert McCann, Wilfrid Gangbo, Alessio Porretta, Dejan Slepcev, Lorenzo Zambotti}
% on various evolutionary models related to Optimal Transport.



\newpage
% (sticky particle, 4th order equations,
% nonlinear mobility) related to Optimal Transport.

% \npind, Yann Brenier, Jose
%   Carrillo, Lorenzo Zambotti}

% \mysec*{EDITORIAL BOARDS}

% \npind
% 2013--\present: Potential Analysis.

% \mysec*{RESEARCH PROJECTS}

% Leader of the Pavia Research unit of the 
% following Projects of National Interest (PRIN)
% funded by the Italian ministry of
% education (MIUR):

% \npind
% 2007-2009: \emph{Optimal Transport and evolution variational
%   problems}. National leader: Luigi Ambrosio

% \npind
% 2010-2012: \emph{Variational, functional-analytic, and optimal
%   transport methods for dissipative evolutions and stability
%   problems}. National leader: Luigi Ambrosio

% \npind
% 2013-2016: \emph{Calculus of Variations}. 
% National leader: Gianni Dal Maso.




\renewcommand{\textsf}[1]{[#1]}
\subsection*{Publications}
\renewcommand{\thefootnote}{\fnsymbol{footnote}}
%\begin{thebibliography}{PSCF05b}
Author of 110 publications, 
with 15000+ citations, H-index 49 according to 
\href{https://scholar.google.com/citations?hl=en&user=w7qH1YYAAAAJ&view_op=list_works&sortby=pubdate}{Google Scholar}.\\

% \href{https://mathscinet.ams.org/mathscinet/}{MSC database}
% attributes to me 107 publications, 
% 6500+ citations by 3400+
% authors, H-index 36. %  (44 publications in the last 10 years, which
% received 1440 citations)
% Seven of the papers listed below belong to the
% group of
% \href{https://images.webofknowledge.com/images/help/WOS/hs_citation_applications.html}{Highly
%   cited papers} according to WOS.
% as of May/June 2022,
% the paper received enough citations to place it in the top 1\%
% of the academic field of Mathematics based on a highly cited threshold for the field and publication year.

% The most relevant contributions of my research activity over the past
% ten years concern:

%\mybreak

% \textbf{Foundation of the theory of metric-measure spaces with
%   \emph{Riemannian Ricci Curvature bounded from below} (the so-called
%   RCD$(K,N)$ condition):}


\begin{itemize}[wide,
  topsep=-0.5\parsep,
  itemsep=-4pt,
  leftmargin=12pt,
  %listparindent=-24pt,
  %itemindent=-12pt
  ]
\renewcommand{\labelitemi}{$\diamond$}
%\bibitem[1]{} %AGS14a]{AGS14a}
%G.~Savar{\'e}.
%\bibitem[3]{}%{AGS14b}


%   \item%
% \href{https://doi.org/10.1112/PLMS/PDV047}{\emph{Convergence of pointed non-compact metric measure
%     spaces and stability of Ricci curvature bounds and heat flows.}}
%   (with A.~Mondino, N.~Gigli)\\
%   \newblock{Proc. Lond. Math. Soc.} 111 (2015): 1071-–1129. Cit.~69
\item
  \textbf{Foundation of the theory of metric-measure spaces with
  \emph{Riemannian Ricci curvature bounded from below} (the so-called
  RCD$(K,N)$ condition):}

\begin{itemize}[
  topsep=-2\parsep,
  itemsep=-2pt,
  leftmargin=12pt,
  % listparindent=-12pt,
  itemindent=-12pt
  ]
\renewcommand{\labelitemii}{}
  
\item%[\bf 1.] 
\href{https://doi.org/10.1007/s00222-013-0456-1}{\emph{Calculus and heat flow in metric measure spaces and applications to
  spaces with {R}icci bounds from below}} (with L.~Ambrosio, N.~Gigli).
\newblock {Invent. Math.} 195 (2014), 289--391. %Cit.~202
%(\footnotemark[1]{})
\\
{\small\textsf{Heat flow and Cheeger energy
in metric--measure spaces: $L^2$ and Optimal Transport theory}}


%\bibitem[2]{}%{AGS14b}
%and G.~Savar{\'e}.
\item%[\bf \bf 2.]
\href{https://doi.org/10.1215/00127094-2681605}{\emph{Metric measure spaces with {R}iemannian {R}icci curvature bounded
  from below}}
  (with L.~Ambrosio, N.~Gigli).
\newblock {Duke Math. J.}, 163
(2014):1405--1490. %Cit.~238 
%(\footnotemark[1]{}).
\\
{\small\textsf{The first analysis and characterization
  of Riemannian RCD(K,$\infty$) spaces with quadratic Cheeger energy}}

\item%[\bf 3.] %, and G.~Savar{\'e}.
  \newblock
  \href{https://doi.org/10.1214/14-AOP907}{\emph{Bakry-\'{E}mery curvature-dimension condition and
{R}iemannian {R}icci curvature bounds}}
(with L.~Ambrosio, N.~Gigli).
\newblock {Annals of Probability},  43 
(2015): 339--404.  %Cit.~122 
%(\footnotemark[1]{})
\\
{\small\textsf{The full identification between the Bakry-\'Emery
  condition and RCD(K,$\infty$) spaces}}

\item%
\href{https://doi.org/10.3934/DCDS.2014.34.1641}{\emph{Self-improvement of the Bakry-\'Emery condition and
    Wasserstein contraction of the heat flow in RCD$(K,\infty)$ metric
    measure spaces}}
\newblock{Discrete Contin. Dyn. Syst.} 34 (2014): 1641–1661. %Cit.~70
%(\footnotemark[1]{})
\\
{\small\textsf{Second order calculus, measure-valued $\Gamma_2$-tensor, and
    improved
    Bakry-\'Emery
  condition}}

\item\href{https://doi.org/10.1112/plms/pdv047}{\em 
  Convergence of pointed non-compact metric measure spaces and
  stability of Ricci curvature bounds and heat flows}
(with N.~Gigli, A.~Mondino)
\newblock {Proc. Lond. Math. Soc.} 111 (2015): 1071–1129. %Cit.~83
%(\footnotemark[1]{})

  \item%
\href{https://doi.org/10.1090/memo/1270}{\emph{Nonlinear diffusion equations and curvature conditions
    in metric measure spaces}}
  (with L.~Ambrosio, A.~Mondino)
  \newblock{Mem. Amer. Math. Soc.} 262 (2019), no. 1270, v+121 pp.
  % ISBN: 978-1-4704-3913-2.
  %Cit.~49
  \\
    {\small\textsf{The full identification between the Bakry-\'Emery
      condition and the RCD$(K,N)$ spaces}}
\end{itemize}

\item
  \textbf{Metric-Sobolev spaces: identification of the construction by
    Lipschitz functions and Cheeger energy with the Newtonian approach}
\begin{itemize}[
  topsep=-2\parsep,
  itemsep=-2pt,
  leftmargin=12pt,
  % listparindent=-12pt,
  itemindent=-12pt
  ]
  \renewcommand{\labelitemii}{}

  \item%
\href{https://doi.org/10.4171/RMI/746}{\emph{Density of Lipschitz functions and equivalence of weak
    gradients in metric measure spaces}}
  (with L.~Ambrosio, N.~Gigli)
  \newblock {Rev. Mat. Iberoam.} 29 (2013): 969--996. %Cit.~96 
  %(\footnotemark[1]{})
\end{itemize}
\item
  \textbf{Entropy-Transport formulation for unbalanced optimal
    transport and the characterization of the new Hellinger-Kantorovich distance}

\begin{itemize}[
  topsep=-2\parsep,
  itemsep=-2pt,
  leftmargin=12pt,
  % listparindent=-12pt,
  itemindent=-12pt
  ]
\renewcommand{\labelitemii}{}

  \item%
\href{https://doi.org/10.1007/S00222-017-0759-8}{\emph{Optimal entropy-transport problems and a new
    Hellinger-Kantorovich distance between positive measures}}
  (with M.~Liero, A.~Mielke).
  \newblock{Invent. Math.} 211 (2018): 969--1117. %Cit.~34 
  %(\footnotemark[1]{})
\end{itemize}

  \item \textbf{Foundation of Balanced Viscosity and Visco-Energetic solutions
to rate-independent processes in infinite-dimensional spaces}

\begin{itemize}[
  topsep=-2\parsep,
  itemsep=-2pt,
  leftmargin=12pt,
  % listparindent=-12pt,
  itemindent=-12pt
  ]
\renewcommand{\labelitemii}{}
  \item%[\bf 4.] %and G.~{Savar{\'e}}.
\href{https://doi.org/10.4171/JEMS/639}{\emph{Balanced viscosity (BV) solutions to infinite-dimensional
  rate-independent systems}} (with A.~Mielke, R.~Rossi).
\newblock {J. Eur. Math. Soc. (JEMS)} 18 (2016): 2107--2165. %Cit.~35
{\small\textsf{The first contribution to existence, characterization,
  and properties of Balanced 
  Viscosity solutions
in infinite dimension.}}

  
% %\bibitem[MRS13a]{Mielke-Rossi-Savare-preprint13}
% \item%[\bf 4.] %and G.~{Savar{\'e}}.
% \newblock {\em Balanced viscosity (BV) solutions to infinite-dimensional
%   rate-independent systems} (with A.~Mielke, R.~Rossi).
% \newblock {JEMS, to appear}; ArXiv 1309.6291:1--41.

%\item 
% \newblock {\em On the duality between $p$-Modulus and probability
%   measures}
% (with L.~Ambrosio, S.~Di Marino).
% \newblock {JEMS, to appear};
% ArXiv e-prints 1311.1381: 1-40.
% {\small
% \textsf{Equivalence of the metric definition of Sobolev spaces
%     based on $p$-modulus or on probability measures on
%     curves}}

% Brenier, Y.; Gangbo, W.; Savaré, G.; Westdickenberg, M. Sticky
% particle dynamics with interactions. J. Math. Pures Appl. (9) 99
% (2013), no. 5, 577–617. (Reviewer: Luisa Arlotti)


% Orrieri, Carlo; Porretta, Alessio; Savaré, Giuseppe A variational
% approach to the mean field planning problem. J. Funct. Anal. 277
% (2019), no. 6,

%  \href{https://doi.org/10.1016/j.matpur.2018.06.022}{\em
%    Weighted energy-dissipation principle for gradient flows in
%  metric spaces.} (with R.~Rossi, A.~Segatti, U.~Stefanelli)
% J. Math. Pures Appl. (9) 127 (2019), 1–66. 

\item%
\href{https://mathscinet.ams.org/leavingmsn?url=https://doi.org/10.1007/s00205-017-1165-5}{\em
  Viscous corrections of the time incremental minimization scheme and
  visco-energetic solutions to rate-independent evolution problems.}
(with L.~Minotti) Arch. Ration. Mech. Anal. 227 (2018), no. 2,
477-543.

\end{itemize}

\item \textbf{The mean-field formulation of spatially inhomogeneous Evolutionary
  Games}

\begin{itemize}[
  topsep=-2\parsep,
  itemsep=-2pt,
  leftmargin=12pt,
  % listparindent=-12pt,
  itemindent=-12pt
  ]
\renewcommand{\labelitemii}{}
\item%
  \href{https://doi.org/10.1002/cpa.21995}{\emph{Spatially Inhomogeneous Evolutionary Games}}
  (with L.~Ambrosio, M.~Fornasier, M.~Morandotti)\\
  \newblock{Comm. Pure Appl. Math.} 74 (2021): 1353--1402
\end{itemize}
\end{itemize}

% In addition to the above papers and topics, in the last ten years I have also been working on
% \begin{smallquote}
%   \emph{sub-linear diffusion, chemical reactions as $\Gamma$-limit of
%     diffusion, Cahn-Hilliard and thin-film equations with nonlinear
%     mobility,     viscoplasticity at finite strain,
%     sticky particle dynamics with
%     interactions,
%     mean-field optimal control and mean-field planning,
%     the concavity of R\'enyi entropy power,
%     variational convergence of gradient flows and
%     rate-independent evolutions,
%     the Weighted Energy-Dissipation
%     principle for gradient flows, reverse
%     approximation of gradient flows, jump processes as generalized
%     gradient flows,
%     the Opial property in Wasserstein
%     spaces, duality
%     properties of metric Sobolev spaces and capacity.}
% \end{smallquote}

%\footnotetext[1]{\emph{A Highly cited paper according to WOS}}

%\newpage

\subsection*{Monographs and contributions to volumes}

\begin{itemize}[topsep=-0.5\parsep,itemsep=-4pt,leftmargin=12pt,listparindent=-12pt,itemindent=-12pt] 
\renewcommand{\labelitemi}{}
%\bibitem[AGS08]{MR2401600}%
%L.~Ambrosio, N.~Gigli, and G.~Savar{\'e}.
\item
\href{https://doi.org/10.1007/978-3-030-84141-6}{\em Sobolev Spaces in Extended Metric-Measure Spaces} 
\newblock in \emph{New Trends on Analysis and Geometry in Metric Spaces.}
Lecture Notes in Mathematics, Springer 2022, 117--276.\\
{\small\textsf{The refined theory of metric Sobolev space generated by
    sub-algebra of Lipschitz functions in extended metric-measure spaces.}}
    

\item \href{https://link.springer.com/book/10.1007/b137080}{\em Gradient flows in metric spaces and in the space of probability
   measures} (with L.~Ambrosio, N.~Gigli)
\newblock Lectures in Mathematics ETH Z\"urich. Birkh\"auser Verlag, Basel,
 2005 (second edition 2008).\\
 The two editions have received more than
1500 
 citations, according to MSC.

\item
\href{https://link.springer.com/chapter/10.1007/88-470-0396-2_6}{\em 
Computational electrocardiology: mathematical and numerical modeling}
 (with P.~Colli~Franzone, L.~F. Pavarino; contribution).
 \newblock In {Complex systems in biomedicine}, pages 187--241. Springer
 Italia, Milan, 2006.







\end{itemize}
\subsection*{Invited presentations to conferences (selection)}
\newcommand{\noemph}[1]{#1}
\begin{itemize}[topsep=-0.5\parsep,itemsep=-4pt,leftmargin=*] 
  \renewcommand{\labelitemi}{--}
  \item mSPACE Kickoff Meeting:
  \href{https://bidsa.unibocconi.eu/multiscale-stochastics-patterns-and-analysis-combinatorial-environments}{\em Multiscale Stochastics, Patterns, and Analysis of Combinatorial Environments}, Bocconi University, Milano, 2026
  \item CIRM Workshop:
  \href{https://conferences.cirm-math.fr/3049.html}{\em Aggregation-Diffusion Equations \&\ Collective Behavior:
Analysis, Numerics and Applications}, Marseille, 2024
\item 
  MFO Workshop: 
\href{https://www.mfo.de/occasion/2406/www_view}{\em Applications of
    Optimal Transportation}, Oberwolfach, {2024}
  \item 
  \workshop{ICMS Workshop:   \href{https://www.icms.org.uk/workshops/2023/optimal-transport-and-calculus-variations}{Optimal Transport and the Calculus of Variations}}{Edinburgh}{2023}{Geodesic convexity of entropy functionals in Hellinger-Kantorovich metric}
  \item 
  \workshop{Workshop on 
  \href{https://www.msp2023.campus.unimib.it}{The Mathematics of Subjective Probability}}{Milano}{2023}{Sobolev spaces on the Wasserstein space of probability measures}
  \item 
  \workshop{Workshop \href{https://events.dm.unipi.it/event/77}{Calculus of Variations \& Geometric Measure Theory}}{Pisa}{2023}{Evolution of probability measures: beyond gradient flows}
  \item 
  \workshop{BIRS Workshop: \href{https://www.birs.ca/events/2023/5-day-workshops/23w6003}{Nonlinear Diffusion and nonlocal Interaction Models - Entropies, Complexity, and Multi-Scale Structures}}{Granada}{2023}{A Lagrangian approach to dissipative evolutions of probability measures}
  \item \workshop{Workshop \href{https://brinmrc.umd.edu/programs/workshops/spring23/spring23-workshop-numerics.html}{\noemph{Frontiers of Numerical PDEs}}}{Maryland}{2023}{Evolution equations in spaces of probability measures}
  \item \workshop{Workshop \href{https://sites.google.com/view/ot-lagrange/barycenters-interpolation-harmonic-maps?authuser=0}{\noemph{Interpolation of Measures}}}{Paris}{2023}{Fine properties of geodesics and geodesic $\lambda$-convexity for the Hellinger-Kantorovich distance}
\item \workshop{MFO workshop: \href{https://www.mfo.de/occasion/2244/www_view}{\noemph{Heat Kernels, Stochastic Processes and Functional Inequalities}}}{Oberwolfach}{2022}{Capacitary modulus and Newtonian-Sobolev capacity in metric measure spaces}
\item \workshop{HCM Conference: \href{https://www.hcm.uni-bonn.de/events/eventpages/2022/from-dirichlet-forms-to-wasserstein-geometry/}{\noemph{From Dirichlet Forms to Wasserstein
    Geometry}}}{Bonn}{2022}{Density of subalgebras of Lipschitz functions in metric Sobolev spaces and applications to measured Wasserstein spaces}
\item \workshop{Workshop on \href{https://people.maths.ox.ac.uk/carrillo/Capri/anacapri4.html}{Frontiers in Nonlocal Nonlinear PDEs}}{Anacapri}{2022}{Dissipative evolutions of probability measures}  \item \workshop{Workshop \href{https://sites.google.com/uniroma1.it/summerschool2022/home-page?authuser=0}{\noemph{20 years of Summer Schools on CalcVar in Rome}}}{Roma}{2022}{Evolution of probability measures}
% \item Workshop on \href{http://people.maths.ox.ac.uk/carrillo/Capri/anacapri4.html}{\em Frontiers in Nonlocal Nonlinear PDEs}, Anacapri, 2022
\item \workshop{Workshop \href{https://sites.google.com/view/gmtcortona2020/?pli=1}{Geometric Measure Theory and applications}}{Cortona}{2021}{Lipschitz approximation, Capacity and Capacitary Modulus in metric Newtonian-Sobolev spaces}
\item \workshop{MFO workshop: \href{https://publications.mfo.de/handle/mfo/3858}{\noemph{Applications of Optimal Transportation in the Natural Sciences}} (online)}{Oberwolfach}{2021}{Dissipative evolution of measures}
\item \workshop{Workshop: \href{https://gianni65.weebly.com}{\noemph{Calculus of Variations and Applications}}}{SISSA, Trieste}{2020}{Singular perturbation of gradient flows and rate-independent evolution}
\item \href{https://umi.dm.unibo.it/congresso2019/}{\noemph{XXI Congresso Unione Matematica Italiana}}, 2019. {Plenary Speaker}
\item Workshop: \href{https://otgeoan.wixsite.com/venice}{\noemph{Optimal transport and Geometric Analysis}}, Venice, 2019.
\item ICMS Workshop: \href{https://www.icms.org.uk/workshops/2018/gradient-flows-challenges-and-new-directions}{\noemph{Gradient flows: challenges and new directions}, Edinburgh}, 2018.
\item BIRS workshop:
  \href{https://www.birs.ca/events/2018/5-day-workshops/18w5094}{\noemph{Topics in the Calculus of Variations: Recent Advances and New Trends}}, Banff, 2018.
\item BIRS workshop: \href{https://www.birs.ca/events/2018/5-day-workshops/18w5069}{\noemph{Entropies, the Geometry of Nonlinear Flows, and their Applications}}, Banff, 2018.
\item MFO workshop: \href{https://publications.mfo.de/handle/mfo/3618}{\noemph{Variational Methods for Evolution} Oberwolfach}, 2017
\item MFO workshop: \href{https://publications.mfo.de/handle/mfo/3571}{\noemph{Applications of Optimal Transportation in
    the Natural Sciences} Oberwolfach}, 2017
\item  7$^{th}$ European Congress of Mathematics,
Berlin: section ``Analysis and PDEs''. Invited speaker.
\item MFO workshop: \href{https://publications.mfo.de/handle/mfo/3559}{\noemph{Heat Kernels, Stochastic Processes, and Functional Inequalities}}, Oberwolfach, 2016.
\item Conference on \href{https://wt.iam.uni-bonn.de/conference-new-trends-in-optimal-transport}{\noemph{New trends in Optimal Transport}}, HIM, Bonn, 2015.
\item BIRS workshop: \href{https://www.birs.ca/events/2014/5-day-workshops/14w5109}{\noemph{Entropy Methods, PDEs, Functional Inequalities, and
  Applications}},
Banff, 2014.
\item International Conference on
\noemph{Fractal Geometry and Stochastics V}, Tabarz, 2014.
Plenary
speaker.

\item Workshop on
  \href{https://www.msri.org/workshops/728}{\noemph{Infinite-Dimensional Geometry}},
MSRI, Berkeley, 2013.


\item EQUADIFF 2013, Prague. Plenary speaker.

% \item Conference on \noemph{Probability and Geometry},
% Poitiers, 2012.

\item BIRS Workshop: \noemph{Optimal Transportation and Differential Geometry}
Banff, 2012.

\item MFO workshop: \href{https://publications.mfo.de/handle/mfo/3277}{\noemph{Interplay of Analysis and Probability in Physics}},
Oberwolfach, 2012.

\item MFO mini-workshop: \href{https://publications.mfo.de/handle/mfo/3273}{\noemph{Manifolds with Lower Curvature Bounds}},
Oberwolfach, 2012.

\item RISM meeting: \href{https://www.mate.polimi.it/workshop/RISM3/index.php?tab=home}{\noemph{Multiphase and Multiphysics problems}}, 
 Verbania (IT), 2011.

 \item BIRS workshop: \href{http://www.birs.ca/events/2010/5-day-workshops/10w5054}{\noemph{Nonlinear Diffusions and Entropy
   Dissipation: From Geometry to Biology}}, %Banff, 
   2010.

 \item CIRM-HCM Meeting:
 \noemph{Stochastic Analysis, SPDEs, Particle Systems, Optimal Transport}
 %Levico (IT), 
 2010.

\item Workshop on \href{https://www.him.uni-bonn.de/programs/past-programs/past-trimester-programs/complex-stochastic-systems-discrete-vs-continuous/workshop-3-particle-systems-nonlinear-diffusions-and-equilibration/}{\noemph{Particle systems, nonlinear diffusions, and equilibration}},
 HCM, Bonn, 2007.

 \item Workshop on 
 \noemph{Optimal Transportation, and Applications to Geophysics and Geometry},
 Edinburgh, 2007.

\item Workshop on \noemph{Optimal transport: theory and application}, Centro De Giorgi, Pisa, 2006.

% \item \noemph{Nonlinear Diffusion Equations and related PDEs},
%  UAM, Madrid, 2006

 % \item Workshop on \noemph{Modelling and analysis of phase transitions}
 % Centro De Giorgi, Pisa, 2006.

 \item ICMS Workshop: \noemph{Optimal Transportation, Transport Equations and Hydrodynamics},
 Edinburgh, 2005.

 \item \noemph{10th Conference on Free Boundary Problems},
 Coimbra, June 7-12, 2005. Plenary speaker.
\end{itemize}

\subsection*{Invited courses to international advanced schools}

\begin{itemize}[topsep=-0.5\parsep,itemsep=-4pt,leftmargin=*] 
  \renewcommand{\labelitemi}{--}

\item Spring School
\href{https://www.uni-ulm.de/mawi/iaa/forschung/spring-school-2026-horizons-in-nonlinear-pdes/}{\emph{Horizons in Nonlinear PDEs}},
  Universit\"at Ulm, Germany, 2026:
  Variational Principles for Evolution Problems

\item CIME course on
\href{https://link.springer.com/book/10.1007/978-3-030-84141-6}{\emph{New Trends on Analysis and Geometry in
    Metric Spaces}},
  Levico Terme, Italy 2017:
  Sobolev Spaces in Extended Metric-Measure Spaces
  
\item
\href{https://www.him.uni-bonn.de/programs/past-programs/past-junior-trimester-programs/optimal-transportation/workshop-gradient-flows-and-entropy-methods/}{\emph{Gradient flows and entropy methods}}, HIM, Bonn, 2015:
The Weighted Energy-Dissipation (WED) principle for gradient flows.


\item 
\href{https://cvgmt.sns.it/event/270/}{\emph{Analysis and Geometry on Singular Spaces}},
Scuola Normale Superiore, Pisa, 2014:
Metric measure spaces with Riemannian Ricci curvature bounded from below.

\item
\href{https://www1.mat.uniroma1.it/people/garroni/School2013.html}{\emph{Seventh Summer School in Analysis and Applied Mathematics}},
Roma, 2013: Gradient flows and rate-independent evolutions: a variational approach.

 \item
CNA Summer School on
\href{https://www.math.cmu.edu/cna/2010CNASummerSchoolFiles/index.html}{\emph{ "New Vistas in Image Processing and PDEs"}}
Carnegie Mellon University, Pittsburgh, 2010:
Applications of optimal transport to evolutionary PDEs.

\item 
 \emph{School on "Optimal transport: Theory and applications"}
 Institut Fourier, Grenoble, 2009: Gradient flows and optimal transport.

 \item 
 \href{https://www2.karlin.mff.cuni.cz/~prusv/ncmm/conference/eveq2008/info.html}{\emph{EVEQ2008}},
 Prague, 2008: A variational approach to gradient flows and rate-independent problems.

 \item 
 \emph{School in Nonlinear Analysis and Calculus of Variations}
 Scuola Normale Superiore, Pisa, 2006:
 Gradient flows: a variational approach.
\end{itemize}

% \medskip
% DICHIARAZIONE SOSTITUTIVA DI CERTIFICAZIONE (ART. 46 E 47 D.P.R. 445/2000)
% Il sottoscritto Giuseppe Savar\'e  ai sensi e per gli effetti degli articoli 46 e 47 e consapevole delle sanzioni penali previste dall’articolo 76 del D.P.R. 28 dicembre 2000, n. 445 nelle ipotesi di falsità in atti e dichiarazioni mendaci, dichiara che le informazioni riportate nel presente curriculum vitae corrispondono a verità.

\vspace{.5truecm}
Milano, February 14, 2026

\begin{center}
  Giuseppe Savar\'e
\end{center}
\newpage
% \begin{center}
% Giuseppe Savar\'e - List of publications
% \end{center}

\nocite{*}
{\renewcommand{\markboth}[2]{}% Remove header adjustment
\printbibliography[title={List of publications}]}


\end{document}